\section{Without Which No Man Shall See God}

The promise of the vision carries its condition inside it. ``Blessed are the clean of heart: for they shall see God.'' The Epistle to the Hebrews states the same condition from the other side: ``Follow peace with all men and holiness: without which no man shall see God.''\endnote{Matt 5:8; Heb 12:14, Douay--Rheims (Challoner).} And the Apocalypse, having opened the gates of the city that are never shut, says who will pass through them: ``There shall not enter into it any thing defiled or that worketh abomination or maketh a lie: but they that are written in the book of life of the Lamb.''\endnote{Apoc 21:27, Douay--Rheims (Challoner).} Those who die in God's grace but not yet perfectly purified are purified after death, ``so as to achieve the holiness necessary to enter the joy of God'';\endnote{\emph{Catechism of the Catholic Church} 1054.} that mercy shown to the saved is the subject of \emph{What Death Does Not Finish}. Holiness and heaven belong together for a reason.

\subsection{Heaven would be hell to an irreligious man}

John Henry Newman gave the answer in the first of his \emph{Parochial and Plain Sermons}, preached on that text from Hebrews. He asks why we cannot be saved without holiness, and replies that an unholy man, even if he were let into heaven, could not be happy there. Heaven, he says, is not a place where everyone takes his own pleasure, as on earth:

\begin{quote}
Heaven then is not like this world; I will say what it is much more like,---a church. For in a place of public worship no language of this world is heard; there are no schemes brought forward for temporal objects, great or small \ldots\ Here we hear solely and entirely of God. We praise Him, worship Him, sing to Him, thank Him, confess to Him, give ourselves up to Him, and ask His blessing.\endnote{John Henry Newman, ``Holiness Necessary for Future Blessedness,'' \emph{Parochial and Plain Sermons}, vol.~1, sermon~1 (London: Rivingtons, 1869).}
\end{quote}

A man who has never loved God, brought into that presence, would find in it nothing he cared for and everything he had avoided. ``Ah! he could not bear the face of the Living God; the Holy God would be no object of joy to him. \ldots\ None but the holy can look upon the Holy One; without holiness no man can endure to see the Lord.'' And then he goes further:

\begin{quote}
Nay, I will venture to say more than this;---it is fearful, but it is right to say it;---that if we wished to imagine a punishment for an unholy, reprobate soul, we perhaps could not fancy a greater than to summon it to heaven. Heaven would be hell to an irreligious man.\endnote{Newman, \emph{Parochial and Plain Sermons}, vol.~1, sermon~1.}
\end{quote}

``Fire does not inflame iron, but it inflames straw,'' Newman says. The same holiness that is joy and life to the saint would be torment to the soul that hates it. And so holiness is not an arbitrary fee for admission. It is the capacity to enjoy what heaven is. ``Outward acts, done on principle, create inward habits,'' he adds; every prayer, every act of charity and self-denial, schools the heart for the presence of God.\endnote{Newman, \emph{Parochial and Plain Sermons}, vol.~1, sermon~1.} Thomas said the same in his own terms when he taught that charity measures the light of glory. The heart made clean by grace is the eye made ready to see.

\subsection{Meditating on paradise}

The saints did not leave this at doctrine. They taught the faithful to meditate on heaven as they meditated on death, judgment and hell, and to draw from the meditation affections and resolutions. Francis de Sales, in the eighth meditation of the \emph{Introduction to the Devout Life}, bids his reader place himself in the presence of God and then consider:

\begin{quote}
Represent to yourself a lovely calm night, when the heavens are bright with innumerable stars: add to the beauty of such a night the utmost beauty of a glorious summer's day,---the sun's brightness not hindering the clear shining of moon or stars, and then be sure that it all falls immeasurably short of the glory of Paradise.\endnote{Francis de Sales, \emph{Introduction to the Devout Life} I.16, eighth meditation, ``On Paradise'' (London: Rivingtons, 1876).}
\end{quote}

Then the company of angels and saints, and then the chief joy: ``Consider how they enjoy the Presence of God, Who fills them with the richness of His Vision, which is a perfect ocean of delight.'' From the considerations he draws the affections, and the second of them is a reproach to the heart: ``Reprove the coldness of your own heart for having hitherto so little sought after that glorious abode. Why have I so long lingered indifferent to the eternal happiness set before me?'' And the resolution is practical, as a saint's resolution always is: ``Resolve to give up such and such things, which hinder you on the way, and to do such others as will help you thitherwards.''\endnote{Francis de Sales, \emph{Introduction to the Devout Life} I.16, considerations 3 and affections 2 and 4.}

Alphonsus ends his consideration on heaven by turning the hope into patience. ``When, therefore, the crosses of this life afflict us, let us animate ourselves with the hope of Heaven to bear them patiently.'' He tells of Philip Neri, who, when the dignity of cardinal was offered him, ``threw up the cap in the air, exclaiming, Paradise! Paradise!''; and he gives his reader the same short prayer for every trouble: ``Let us likewise, when we are afflicted by the miseries of this life, raise up our eyes to Heaven, and console ourselves, saying with a sigh, Heaven! Heaven!'' For ``if we be faithful to God, all these sorrows, miseries, and fears, shall one day have an end, and we shall be admitted into that blessed country, where we shall enjoy complete happiness as long as God shall be God.''\endnote{Alphonsus, \emph{Preparation for Death}, Consideration XXIX, third point (Duffy, 1843).}

\subsection{The resolution}

The season gives the resolution its shape. The year ends; the Church prays for her dead and looks for her Lord; Advent will soon bid her watch. To desire heaven in such a season is to do three things.

It is to keep the heart clean, by confession and by guarding the eyes and the will, because only the clean of heart shall see God. It is to grow in charity, in the small and hidden acts that widen the soul, because charity is the measure of the light by which God will be seen. And it is to keep desire alive, not quieting the restless heart with lesser goods, but turning its restlessness toward the face it was made for. ``And every one that hath this hope in him sanctifieth himself, as he also is holy.''\endnote{1 John 3:3, Douay--Rheims (Challoner).}

The prayer for such a resolution is one the saints have already written. Alphonsus closes his meditation with an affection that any Christian may make his own:

\begin{quote}
Heaven! Heaven! When, O Lord, shall the day come, on which I shall see you face to face, and embrace you, without fear of ever more losing you? Ah my God, keep your hand upon me, that I may never more offend you.\endnote{Alphonsus, \emph{Preparation for Death}, Consideration XXIX, third point, affections and prayers (Duffy, 1843).}
\end{quote}
